\section{Related Work}
\label{sec:related_work}

Our work is positioned at the intersection of parameter-efficient model adaptation, weight-space merging, dynamic ensembling, and the application of physical and biological principles to deep learning architectures.

\subsection{Parameter-Efficient Fine-Tuning and Model Merging}
Parameter-Efficient Fine-Tuning (PEFT) techniques, particularly Low-Rank Adaptation (LoRA) \cite{hu21lora}, have revolutionized how foundational models are adapted. Instead of updating all model parameters, PEFT trains a tiny fraction of parameters (e.g., adapters \cite{houlsby19adapter} or prefix tuning \cite{liliang21prefix}) while keeping the base network frozen. When multiple task-specific adapters are available, \emph{model merging} aims to combine their capabilities into a single model without incurring additional inference costs. Static model merging techniques like Task Arithmetic \cite{ilharco22taskarithmetic}, TIES-Merging \cite{yadav23ties}, and DARE \cite{yu23dare} resolve parameter interference and average weight updates. However, these methods produce a single set of static weights. When subjected to heterogeneous streams where distinct samples require different expertise, static merged models experience severe interference and performance degradation, termed \emph{Heterogeneity Collapse}.

\subsection{Dynamic Ensembling and Adaptive Routing}
To overcome the limitations of static merging, dynamic ensembling and mixture-of-experts (MoE) \cite{shazeer17moe} have been proposed. MoE architectures route inputs to specialized networks dynamically using learned parametric gating networks. While highly effective, standard MoEs (e.g., Switch Transformers \cite{fedus21switch} or Vision MoEs \cite{riquelme21vmoe}) require expensive joint training of the routers and experts from scratch, making them unsuitable for post-hoc ensembling of independently trained models. Furthermore, parametric routers are highly susceptible to \emph{Vectorization Collapse} in sample-wise streaming (batch size $B=1$), where unregularized parametric routers produce erratic, highly unstable routing decisions when deprived of batch-level statistics. 

\subsection{Streaming Adaptation and Test-Time Merging}
In resource-constrained streaming scenarios, several post-hoc, training-free dynamic merging methods have emerged. SABLE \cite{liang23sable} performs sample-wise activation blending, and SPS-ZCA \cite{zhang24spszca} utilizes zero-shot centroid alignment in early shared layers to route samples. To handle extremely mixed-task streams, systems-level scheduling wrappers like Micro-Batch Homogenization (MBH) buffer incoming samples and group them by estimated task identity to prevent heterogeneity collapse. However, MBH's scheduling introduces high latency and requires $O(K)$ sequential backbone execution passes, which is unacceptable for real-time edge hardware. Crucially, existing post-hoc routing methods (SABLE, SPS-ZCA) treat each layer as an isolated routing block. They evaluate stateless similarities at each block, which leads to sharp, high-frequency ensembling weight oscillations (routing jitter) and representational drift across deep layers.

\subsection{Physical and Dynamical Systems in Deep Learning}
Connecting deep learning to physical systems is a rich and fruitful avenue of research. Classic examples include Hopfield Networks \cite{hopfield1982neural}, which model associative memory via spin-glass energy landscapes. More recently, the interpretation of residual networks (ResNets) as discretizations of ordinary differential equations has led to the development of Neural ODEs \cite{chen2018neuralode} and stable, continuous-depth architectures \cite{haber2017stable, lu2018beyond, grathwohl2018ffjord}. While these works model individual sample representations as trajectories of dynamical systems, they do not consider multi-task model ensembling. In ChemMerge, we leverage this continuous-time perspective to model the ensembling process itself as a non-equilibrium chemical reactor. By treating task expert activations as concentration-dependent reacting species, we introduce physical inertia into the ensembling weight trajectory, resolving the stateless routing jitter paradox.
